%=====================================================================
% PART IV OPENER
%=====================================================================
\thispagestyle{plain}
\structuralfigures{P4.}

\begin{center}
\begin{tcolorbox}[enhanced, width=0.92\textwidth, colback=purplehl!5,
  colframe=purplehl, boxrule=0pt, leftrule=5pt, arc=3pt, top=8pt, bottom=8pt]
\large\itshape\color{purplehl!70!black}
Divide-and-conquer splits a problem into independent pieces. Part~IV is what to do when the pieces are \emph{not} independent --- when a local choice can be trusted, when it cannot and the subproblems overlap, and how to cost a sequence of operations rather than one.
\end{tcolorbox}
\end{center}

\vspace{6pt}

\dsfigh{%
\begin{tikzpicture}[
  cbox/.style={rounded corners=3pt, draw=purplehl, line width=1pt, fill=purplehl!7,
               text=purplehl!70!black, font=\small\bfseries, align=center,
               minimum width=3.5cm, minimum height=1.0cm},
  hbox/.style={rounded corners=3pt, draw=purplehl, line width=1.8pt, fill=purplehl!16,
               text=purplehl!70!black, font=\small\bfseries, align=center,
               minimum width=3.5cm, minimum height=1.0cm},
  arr/.style={-{Stealth[length=2.4mm]}, purplehl!60, line width=1pt},
  note/.style={font=\scriptsize\itshape, text=purplehl!70!black, align=center,
               fill=white, inner sep=2pt}
]
  \node[cbox] (c0) at (0.00,0) {Ch.\ 14\\Greedy};
  \node[hbox] (c1) at (4.30,0) {Ch.\ 15\\DP: Elements};
  \node[cbox] (c2) at (8.60,0) {Ch.\ 16\\DP: Problems};
  \node[cbox] (c3) at (12.90,0) {Ch.\ 17\\Amortised};
  \draw[arr] (c0) -- (c1);
  \draw[arr] (c1) -- (c2);
  \draw[arr] (c2) -- (c3);
  \node[note] at (2.15,-1.0) {when it cannot};
  \node[note] at (6.45,-1.0) {the classics};
  \node[note] at (10.75,-1.0) {costing a sequence};
  \node[font=\scriptsize\itshape, text=accent, align=center] at (4.30,1.15)
       {the two properties that decide which applies};
\end{tikzpicture}}%
{Reading path through Part~IV. Chapter 15 is the one to understand rather than memorise: the two properties in it decide whether greedy or dynamic programming applies at all.}{part4map}

\newpage
